\section{The question, the published obstruction, and what the record claims}\label{sec:problem}

\subsection{The question}

Whether the six-sphere carries a complex structure (an integrable almost complex structure) has long been open. Campana, Demailly and Peternell studied a hypothetical complex structure $X$ on $S^6$ through its algebraic dimension $a(X)$, the transcendence degree of its field of meromorphic functions, first in 1998 \cite{CDP98} and then in 2020 \cite{CDP20}. In the 2020 paper they prove that a compact complex threefold homeomorphic to $S^6$ would have algebraic dimension $0$ (\cite{CDP20}, Corollary~2.3; Engel's exposition cites the same statement from the 1998 paper, \cite{CDP98}, Corollary~4.2). A complex structure on $S^6$ that fibres holomorphically over $\Pp^1$ has $a(X)\ge1$, and would therefore contradict that theorem.

\subsection{The claimed construction, the record, and an exposition}

\begin{itemize}
\item \emph{The source manuscript} \cite{AlpogeS6}, whose opening summary is titled \emph{The $(3,4,\infty)$ modular family of 2-tori, completed at its three special points, is a complex structure on $S^6$}, claims exactly such a threefold: $X\to\Pp^1$ with $a(X)=1$, $X$ diffeomorphic to $S^6$. The record reconstructs it as file \texttt{03\_} (108 pages). It was circulated by Levent Alpöge and produced with Claude, according to the attribution file of the Yang--Mills workbench.
\item \emph{The record} (24 August to 9 September 2026), by KokunoYumeto, audits that manuscript. It reconstructs the manuscript, repairs local defects, and concludes, with its own label \textsf{VERIFIED}, that the constructed threefold satisfies every hypothesis of the relevant theorem of \cite{CDP20} and contradicts its conclusion. It also isolates a step in the published proof that is not a formal consequence of the sequences displayed there. The source manuscript had already located the failure of that step at its own cusp fibre; the record's correction note adds an independent test family (Sections~\ref{sec:ledger} and~\ref{sec:correction}).
\item \emph{An exposition.} Philip Engel's \emph{Complex structures on $S^6$} (13 September 2026) \cite{EngelS6} presents the construction, which it credits to Claude under Alpöge's direction, with the stated aim of a self-contained proof. It notes that the construction contradicts the result that the algebraic dimension must be $0$, which it cites as \cite{CDP98}, Corollary~4.2, and it describes the construction as a one-parameter family of threefolds diffeomorphic to $S^6$.
\end{itemize}
The global construction and the counterexample claim are not verified here. This paper reports what the record claims, and proves or checks what can be established from the record's local, finite and algebraic statements.

\subsection{The record's own ledger and scope}

\begin{itemize}
\item At the freeze of version 1.0 its internal claim ledger (in the evidence archive \texttt{04\_}) had $42$ rows labelled \textsf{VERIFIED}, none \textsf{CONDITIONAL}, $2$ \textsf{OPEN} and $4$ \textsf{REFUTED}. The two open rows are:
\begin{itemize}
\item independent specialist replication of the cyclic invariant lattices, the logarithmic determinant divisor, the residue cancellation and an order-four scalar test;
\item specialist and community adjudication of the complete construction and counterexample.
\end{itemize}
The README describes them as replication and external adjudication, not missing internal premises.
\item The later supplements (5--9 September) use the construction's data as stated input, and state that they do not by themselves establish a complex structure on $S^6$ or disprove the published theorem.
\end{itemize}
